% =====================================================================
% NegoPlay — Business Plan
% Author: Anna Ben-Shushan
% Supervisor: Dr. Yoram Segal
% Course: AI Development Expert
% =====================================================================
\documentclass[11pt, a4paper, oneside]{report}

\usepackage[T1]{fontenc}
\usepackage[utf8]{inputenc}
\usepackage{lmodern}
\usepackage{microtype}
\usepackage[margin=2.5cm]{geometry}
\usepackage{setspace}
\onehalfspacing

\usepackage{graphicx}
\usepackage{booktabs}
\usepackage{array}
\usepackage{tabularx}
\usepackage{multirow}
\usepackage{colortbl}
\usepackage[table]{xcolor}
\usepackage{float}
\usepackage{amsmath}
\usepackage{enumitem}
\usepackage[hidelinks, colorlinks=true, linkcolor=black,
            citecolor=blue, urlcolor=blue]{hyperref}
\usepackage{fancyhdr}
\usepackage{titlesec}

% ----- Colors -----
\definecolor{headbg}{HTML}{2C3E50}
\definecolor{accentblue}{HTML}{2980B9}
\definecolor{accentgreen}{HTML}{27AE60}
\definecolor{accentorange}{HTML}{E67E22}
\definecolor{lightgray}{HTML}{F2F3F4}
\definecolor{darkgray}{HTML}{566573}

% ----- Header / Footer -----
\pagestyle{fancy}
\fancyhf{}
\fancyhead[L]{NegoPlay -- Business Plan}
\fancyhead[R]{\leftmark}
\fancyfoot[C]{\thepage}
\renewcommand{\headrulewidth}{0.4pt}

% ----- Chapter style -----
\titleformat{\chapter}[hang]
  {\normalfont\Huge\bfseries}{\thechapter}{20pt}{\Huge}

% ----- Helper: highlight box -----
\usepackage{tcolorbox}
\tcbuselibrary{skins}
\newtcolorbox{keybox}[1][]{
  colback=accentblue!8, colframe=accentblue,
  fonttitle=\bfseries, title=#1, boxrule=0.8pt
}
\newtcolorbox{metricbox}{
  colback=accentgreen!10, colframe=accentgreen,
  boxrule=0.8pt, fontupper=\bfseries\large\centering
}

% =====================================================================
\begin{document}

% -------------------- Cover Page --------------------
% Removed: bundled as Part III of final_portfolio.tex, which already carries
% the single title page for the whole document plus a "Part III - Business
% Plan" divider. The bundled plan now opens straight at Chapter 1.

% -------------------- TOC --------------------
% No visible TOC here (bundled as Part III of final_portfolio.tex). We still
% force the .toc file open so the master portfolio can read real page numbers;
% \chapter only writes to .toc while it is open (\tableofcontents normally
% opens it), so we open it explicitly without printing anything.
\makeatletter
\@ifundefined{tf@toc}{\newwrite\tf@toc}{}
\immediate\openout\tf@toc\jobname.toc\relax
\makeatother

% =====================================================================
% CHAPTER 1 — EXECUTIVE SUMMARY
% =====================================================================
\chapter{Executive Summary}
\label{ch:exec}

\begin{keybox}[The Business in Three Sentences]
\textbf{NegoPlay} is a B2B SaaS platform that automatically discovers
decision-making profiles from elite competitive data and deploys them as
AI agents to help professionals prepare for, simulate, and win business
negotiations.

The core insight: 149{,}208 expert bridge hands reveal stable
behavioral archetypes (``Slam Hunter'', ``Insurance Player'', etc.)
that transfer measurably to negotiation behavior. We turn this empirical
finding into a subscription API for sales teams, procurement departments,
and M\&A advisors.

\textbf{Target revenue in Year~3:} \$480K ARR.
\textbf{Total budget for MVP:} under \$15 (LLM costs only, excluding developer time).
\end{keybox}

\vspace{0.6cm}

\begin{table}[H]
\centering\small
\caption{Business plan summary.}
\begin{tabular}{ll}
\toprule
\textbf{Category} & \textbf{Value} \\
\midrule
Market (TAM)      & \$4.2B (negotiation intelligence + sales AI) \\
Market (SAM)      & \$420M (mid-market + enterprise, English-speaking) \\
Market (SOM Y3)   & \$480K ARR (realistic 5-year ramp) \\
Business model    & B2B SaaS subscription + usage-based API \\
Primary customer  & Sales directors, procurement managers, M\&A teams \\
Pricing (entry)   & \$49 / month per seat \\
Competitive moat  & Empirically validated profiles from real expert data \\
Technology stack  & Python, Gemini Flash 2.0, scikit-learn, GitHub \\
\bottomrule
\end{tabular}
\end{table}

% =====================================================================
% CHAPTER 2 — THE PROBLEM
% =====================================================================
\chapter{The Problem}
\label{ch:problem}

\section{Negotiations fail because we do not know our counterpart}

Every business negotiation involves incomplete information. We know our
own BATNA (Best Alternative to a Negotiated Agreement) but rarely know
the other side's risk tolerance, decision style, or likely concession
pattern. The result:

\begin{itemize}[leftmargin=*]
    \item Deals are left on the table when one side concedes too quickly.
    \item Deals collapse when one side misjudges the other's walk-away point.
    \item Sales cycles are prolonged because reps cannot adapt their style
          to the buyer's decision profile.
\end{itemize}

\section{The data gap}

Business negotiation lacks the clean feedback loops that make other
domains trainable:

\begin{itemize}[leftmargin=*]
    \item Real negotiation transcripts are proprietary.
    \item Outcomes (``who won?'') are rarely measurable in standardized form.
    \item No public dataset allows training behavioral models on real expert
          negotiation decisions.
\end{itemize}

Existing tools (Gong, Chorus, Clari) analyze \emph{what was said} after
the fact. No tool predicts \emph{how the counterpart will behave} before
the conversation starts.

\begin{keybox}[The gap we fill]
NegoPlay is the first platform to build behavioral profiles from a large,
public, measurable expert decision dataset -- bridge tournament records --
and transfer those profiles to negotiation preparation.
\end{keybox}

% =====================================================================
% CHAPTER 3 — THE SOLUTION
% =====================================================================
\chapter{The Solution}
\label{ch:solution}

\section{What NegoPlay does}

NegoPlay discovers four to five decision-making archetypes from 149{,}208
elite bridge tournament hands. Each archetype is translated into a
named profile with five to seven behavioral traits. The profiles are then
deployed as AI agents that can:

\begin{enumerate}[leftmargin=*]
    \item \textbf{Profile your counterpart} -- given publicly available
          information about the person or company, classify them into the
          nearest archetype.
    \item \textbf{Simulate the negotiation} -- run a multi-round
          AI-vs-AI simulation of the upcoming deal so the user can see
          likely concession patterns before the real meeting.
    \item \textbf{Advise in real time} -- via API integration with CRM
          tools (HubSpot, Salesforce), push profile-based suggestions
          during a live call.
\end{enumerate}

\section{The four profiles (confirmed by Stage 1)}

\begin{table}[H]
\centering\small
\caption{Decision profiles discovered from 149{,}208 bridge tournament hands.
Profile names and assignment rates are the actual Stage~1 outputs (v2.1, May 2026).}
\label{tab:profiles}
\begin{tabularx}{\textwidth}{l c X X}
\toprule
\textbf{Profile} & \textbf{$n$ (\%)} & \textbf{Core traits (from bridge data)} & \textbf{Predicted negotiation behavior} \\
\midrule
Slam Hunter
  & 20 (3.6\,\%)
  & Bids slams at $1.84\times$ the population baseline; calculated upside bets; trusts partner's strength signals.
  & Opens high, concedes little, walks away fast if the value gap is large. \\
Insurance Player
  & 21 (3.7\,\%)
  & Stays in partscore at $1.20\times$ baseline; risk-averse; values certainty over upside.
  & Prefers small guaranteed gains; long close cycles; many contingencies in the contract. \\
NT Specialist
  & 17 (3.0\,\%)
  & Plays NoTrump contracts at $1.36\times$ baseline; analytical; precise trick counting.
  & Data-driven; demands DCF / comparables justification; rejects ``gut-feel'' offers. \\
Fighter
  & 37 (6.6\,\%)
  & Doubles opponents for penalties at $1.55\times$ baseline; competitive; aggressive against opponents.
  & Counter-offers aggressively; threatens walk-away as a tactic; rejects first offer reflexively. \\
\midrule
Generalist (baseline)
  & 468 (83.1\,\%)
  & Average rates on all axes; control population.
  & Default-reasonable; matches counterpart's style. \\
\bottomrule
\end{tabularx}
\end{table}

\noindent\textit{Note on naming consistency.} Earlier internal drafts of this plan used the working names ``Bluffer'' and ``Doubler'' for two of the profiles. After Stage~1 ran on the full 149K-row dataset, the actual statistical axes turned out to be \texttt{nt\_rate} and \texttt{penalty\_double\_rate}, leading to the more descriptive names \emph{NT Specialist} and \emph{Fighter}. The Preliminary Report (Table~A.2) and this Business Plan are now synchronised on these names.

\section{Why bridge data works}

Bridge is one of the only public, large-scale, measurable records of
expert decision-making under incomplete information with four-player
dynamics (two cooperating, two competing -- exactly like a business
negotiation involving two companies each with internal stakeholders).
The IMP scoring system provides a clean, standardized outcome signal.
No comparable dataset exists for business negotiation.

% =====================================================================
% CHAPTER 4 — MARKET ANALYSIS
% =====================================================================
\chapter{Market Analysis}
\label{ch:market}

\section{Total Addressable Market (TAM)}

The global sales intelligence software market was valued at approximately
\$3.4B in 2023 and is growing at 11\,\% CAGR. Adjacent markets include
procurement analytics (\$1.2B) and M\&A advisory tools (\$0.6B). We
define NegoPlay's TAM as the intersection of these segments:

\begin{table}[H]
\centering\small
\caption{TAM breakdown.}
\begin{tabular}{lrr}
\toprule
\textbf{Segment} & \textbf{Market size (2024)} & \textbf{CAGR} \\
\midrule
Sales intelligence (enterprise)    & \$2.8B & 11\,\% \\
Procurement AI tools               & \$0.9B & 15\,\% \\
M\&A deal intelligence             & \$0.5B & 9\,\%  \\
\midrule
\textbf{Total TAM}                 & \textbf{\$4.2B} & \textbf{11\,\%} \\
\bottomrule
\end{tabular}
\end{table}

\section{Serviceable Addressable Market (SAM)}

NegoPlay targets mid-market and enterprise companies in English-speaking
markets (US, UK, Australia, Canada) that: (a) use a CRM, (b) have
active outbound sales or procurement teams, and (c) have an existing AI
software budget. This narrows the addressable market to approximately
10\,\% of the TAM.

\begin{table}[H]
\centering\small
\caption{SAM calculation.}
\begin{tabular}{lr}
\toprule
\textbf{Parameter}                 & \textbf{Value} \\
\midrule
TAM                                & \$4.2B \\
English-speaking market share      & $\times$ 40\,\% \\
Mid-market + enterprise filter     & $\times$ 60\,\% \\
AI-ready companies                 & $\times$ 45\,\% \\
\midrule
\textbf{SAM}                       & \textbf{$\approx$ \$450M} \\
\bottomrule
\end{tabular}
\end{table}

\section{Serviceable Obtainable Market (SOM)}

In the first three years, NegoPlay targets 0.1\,\% of the SAM through
direct sales to early adopters (academic networks, bridge-community
connections, AI-forward SMEs).

\begin{table}[H]
\centering\small
\caption{SOM projection over three years.}
\begin{tabular}{lrrr}
\toprule
\textbf{Metric}               & \textbf{Year 1} & \textbf{Year 2} & \textbf{Year 3} \\
\midrule
Paying customers (seats)      & 20              & 120             & 800 \\
Average revenue per seat/month & \$49           & \$55            & \$50 \\
Monthly Recurring Revenue     & \$980           & \$6{,}600       & \$40{,}000 \\
Annual Recurring Revenue (ARR) & \$11{,}760     & \$79{,}200      & \$480{,}000 \\
\bottomrule
\end{tabular}
\end{table}

% =====================================================================
% CHAPTER 5 — BUSINESS MODEL AND PRICING
% =====================================================================
\chapter{Business Model and Pricing}
\label{ch:bizmodel}

\section{Revenue model}

NegoPlay uses a \textbf{freemium-to-subscription} model with an optional
usage-based API tier for enterprise integrations.

\begin{table}[H]
\centering\small
\caption{Pricing tiers.}
\begin{tabularx}{\textwidth}{l l X r}
\toprule
\textbf{Tier} & \textbf{Audience} & \textbf{Includes} & \textbf{Price} \\
\midrule
Free          & Individual / researcher
              & 10 profile analyses/month, no simulation
              & \$0 \\
Pro           & Sales rep / procurement officer
              & 200 analyses + unlimited simulations + email support
              & \$49/seat/month \\
Team          & SME sales team (up to 20 seats)
              & Unlimited analyses, CRM webhook, team dashboard
              & \$199/month flat \\
Enterprise    & Large enterprise / API integration
              & Unlimited, SLA, SSO, Salesforce / HubSpot connector
              & Custom (est.\ \$2{,}000--\$10{,}000/month) \\
\bottomrule
\end{tabularx}
\end{table}

\section{LLM pricing model -- input vs.\ output tokens}

LLM providers price input tokens (the prompt the model reads) and
output tokens (the response the model generates) at \emph{different}
rates. Output tokens are typically $4\times$ more expensive because
generation requires more compute per token than reading. Any honest
cost model must therefore track both directions separately.

\begin{table}[H]
\centering\small
\caption{Per-million-token pricing across LLM providers (May 2026).
Output is consistently more expensive than input.}
\label{tab:llm_pricing}
\begin{tabular}{l r r r}
\toprule
\textbf{Provider / model} & \textbf{Input} & \textbf{Output} & \textbf{Out/In ratio} \\
\midrule
\textbf{Gemini Flash 2.0 (default)} & \$0.075 / M & \$0.30 / M  & $\mathbf{4.0\times}$ \\
Gemini 2.5 Pro                      & \$1.25 / M  & \$5.00 / M  & $4.0\times$ \\
Claude Haiku 4.5                    & \$1.00 / M  & \$5.00 / M  & $5.0\times$ \\
Claude Opus 4.7                     & \$15.00 / M & \$75.00 / M & $5.0\times$ \\
GPT-4o                              & \$2.50 / M  & \$10.00 / M & $4.0\times$ \\
GPT-4o mini                         & \$0.15 / M  & \$0.60 / M  & $4.0\times$ \\
\bottomrule
\end{tabular}
\end{table}

The cost of a single LLM call therefore decomposes as
\begin{equation}
C_{\text{call}} = \frac{T_{\text{in}}}{10^6}\cdot c_{\text{in}} \;+\; \frac{T_{\text{out}}}{10^6}\cdot c_{\text{out}}
\end{equation}
where $T_{\text{in}}, T_{\text{out}}$ are the input and output token counts
and $c_{\text{in}}, c_{\text{out}}$ are the per-million-token unit prices.

\section{Unit economics per analysis}

The core cost driver is LLM token consumption. Each NegoPlay profile
analysis sends a long prompt (system instruction + counterpart context)
and receives a relatively short structured response. The input/output
breakdown matters because input dominates the token \emph{count} but
output dominates the token \emph{cost}.

\begin{table}[H]
\centering\small
\caption{Unit economics per profile analysis, with input and output costs separated.
Note that 38\,\% of total tokens are output, but they account for 71\,\% of the LLM cost.}
\label{tab:unit_econ}
\begin{tabular}{l r r r}
\toprule
\textbf{Component} & \textbf{Tokens} & \textbf{Unit rate} & \textbf{Cost} \\
\midrule
Gemini Flash 2.0 -- \textbf{input}  (system + user) & 3{,}000 & \$0.075 / M & \$0.000225 \\
Gemini Flash 2.0 -- \textbf{output} (structured JSON) & 1{,}800 & \$0.300 / M & \$0.000540 \\
\midrule
\textbf{LLM subtotal per analysis} & 4{,}800 & --- & \textbf{\$0.000765} \\
Infrastructure (hosting, storage, logging) & --- & --- & \$0.005 \\
\midrule
\textbf{Total cost per analysis}     & --- & --- & \textbf{\$0.0058} \\
\textbf{Revenue per analysis (Pro)}  & --- & --- & \textbf{\$0.245} (200 calls / \$49) \\
\textbf{Gross margin per analysis}   & --- & --- & \textbf{97.6\,\%} \\
\bottomrule
\end{tabular}
\end{table}

\section{Token economics summary across pipeline stages}

\begin{table}[H]
\centering\small
\caption{Token economics across the four pipeline stages, with input
and output costs separated. The Stage 2 skill-extraction call is
the most expensive single call because the LLM reads many bridge
hands (high $T_{\text{in}}$); per-turn negotiation calls are dominated
by output cost.}
\label{tab:token_econ_full}
\begin{tabularx}{\textwidth}{X r r r r r}
\toprule
\textbf{Stage} & \textbf{$T_{in}$} & \textbf{$T_{out}$}
               & \textbf{In cost} & \textbf{Out cost} & \textbf{Total} \\
\midrule
Stage 2 -- Skill extraction       & 7{,}000 & 1{,}000 & \$0.000525 & \$0.000300 & \$0.000825 \\
Stage 3 -- Agent prompt assembly  & 1{,}600 &   400   & \$0.000120 & \$0.000120 & \$0.000240 \\
Stage 4 -- Bridge bidding turn    & 2{,}500 &   500   & \$0.000188 & \$0.000150 & \$0.000338 \\
Stage 4 -- Negotiation turn       & 3{,}500 & 1{,}500 & \$0.000263 & \$0.000450 & \$0.000713 \\
\midrule
\textbf{Full pipeline (one user analysis)} & \textbf{14{,}600} & \textbf{3{,}400} & \textbf{\$0.001096} & \textbf{\$0.001020} & \textbf{\$0.002116} \\
\bottomrule
\end{tabularx}
\end{table}

\begin{keybox}[Why this split matters commercially]
\textbf{18\,\% of NegoPlay's tokens are output, but they account for
48\,\% of the LLM cost.}
Output prompt design is therefore a direct lever on gross margin --
returning concise structured JSON instead of verbose free text cuts the
per-call cost roughly in half. Our agent prompts (see Listing in the
Preliminary Report) require strict JSON schemas of at most 1{,}500
output tokens for this reason.
\end{keybox}

\begin{keybox}[Cost advantage of Gemini Flash 2.0]
With \textbf{input at \$0.075 / M} and \textbf{output at \$0.30 / M},
Gemini Flash 2.0 is approximately $10\times$ cheaper than GPT-4o
(\$2.50 / \$10.00) for equivalent throughput, and roughly $200\times$
cheaper than Claude Opus 4.7 (\$15 / \$75). The entire NegoPlay MVP
(development + testing) runs for under \$15 in LLM costs.
Claude and GPT are used selectively for cross-model validation runs
to strengthen publishability, not for production traffic.
\end{keybox}

% =====================================================================
% CHAPTER 6 — COMPETITIVE LANDSCAPE
% =====================================================================
\chapter{Competitive Landscape}
\label{ch:competition}

\section{Direct and indirect competitors}

\begin{table}[H]
\centering\small
\caption{Competitive landscape.}
\begin{tabularx}{\textwidth}{l X X c}
\toprule
\textbf{Competitor} & \textbf{What they do} & \textbf{Gap vs NegoPlay} & \textbf{Threat} \\
\midrule
\textbf{Gong.io}
  & Analyzes past sales calls (voice + transcript) to coach reps.
  & Post-hoc only; no pre-meeting counterpart profiling; no game-theory basis.
  & Medium \\
\textbf{Pactum AI}
  & Autonomous negotiation bot for procurement (IKEA, Walmart).
  & Fully automated (no human in the loop); no behavioral profiling of individuals.
  & Low--Medium \\
\textbf{Chorus.ai}
  & Records and transcribes sales calls; CI insights for managers.
  & Same as Gong; reactive, not predictive.
  & Low \\
\textbf{Clari}
  & Revenue forecasting + deal risk scoring.
  & Forecasting only; no behavioral simulation.
  & Low \\
\textbf{Crystal Knows}
  & DISC personality profiling from LinkedIn data.
  & Not empirically validated on expert decision data; no simulation.
  & Medium \\
\midrule
\textbf{NegoPlay}
  & Profiles + simulates counterpart from empirically validated archetypes.
  & \emph{Only platform with an empirical game-theoretic foundation.}
  & -- \\
\bottomrule
\end{tabularx}
\end{table}

\section{Our competitive moat}

\begin{enumerate}[leftmargin=*]
    \item \textbf{Empirical foundation:} Profiles are derived from 149{,}208 expert
          decisions with measurable outcomes -- not from personality questionnaires
          or self-reported surveys. This is defensible in peer-reviewed literature.
    \item \textbf{Cross-domain validation:} The alignment finding (Spearman
          $\rho \geq 0.7$ between bridge and negotiation behavior) can be
          published -- creating an academic moat.
    \item \textbf{Provider-neutral architecture:} The unified \texttt{llm\_client.py}
          wrapper means NegoPlay is not locked to any single LLM provider.
    \item \textbf{Data accumulation:} Every new customer interaction refines the
          profiles -- a classic data flywheel.
\end{enumerate}

% =====================================================================
% CHAPTER 7 — GO-TO-MARKET STRATEGY
% =====================================================================
\chapter{Go-to-Market Strategy}
\label{ch:gtm}

\section{Phases}

\subsection*{Phase 1 -- Academic validation (Months 1--6, free)}
\begin{itemize}[leftmargin=*]
    \item Publish the alignment finding as a conference paper (AAMAS or IJCAI).
    \item Open-source the profile discovery code on GitHub to build credibility.
    \item Target bridge clubs, negotiation researchers, and MBA programs as
          early adopters and free validators.
    \item \textbf{Goal:} 100 registered free users, 1 academic citation.
\end{itemize}

\subsection*{Phase 2 -- Paid early access (Months 7--18)}
\begin{itemize}[leftmargin=*]
    \item Launch Pro tier (\$49/seat/month) via Product Hunt and LinkedIn.
    \item Leverage the author's networks at Sami Shamoon College and LUT University.
    \item Offer a 3-month pilot to 3--5 mid-market sales teams at cost.
    \item \textbf{Goal:} 20 paying customers, \$11{,}760 ARR.
\end{itemize}

\subsection*{Phase 3 -- Enterprise and API (Months 19--36)}
\begin{itemize}[leftmargin=*]
    \item Build HubSpot and Salesforce connectors.
    \item Pursue enterprise deals with procurement departments (indirect sales
          via consultancies).
    \item Explore white-labeling for negotiation training companies.
    \item \textbf{Goal:} 800 seats, \$480K ARR.
\end{itemize}

\section{Customer acquisition cost estimate}

Given that Phase 1 is entirely inbound (open source + academic press),
early customer acquisition cost (CAC) is near zero. In Phase 2, assuming
\$2{,}000/month in LinkedIn advertising and content marketing:

\begin{itemize}[leftmargin=*]
    \item \textbf{CAC (Phase 2):} \$2{,}000 / 20 customers $=$ \$100 per customer.
    \item \textbf{LTV (Pro, 12-month retention):} \$49 $\times$ 12 $=$ \$588.
    \item \textbf{LTV/CAC ratio:} 5.9 $\rightarrow$ healthy (target $>$ 3).
\end{itemize}

% =====================================================================
% CHAPTER 8 — ROADMAP AND MILESTONES
% =====================================================================
\chapter{Roadmap and Milestones}
\label{ch:roadmap}

\begin{table}[H]
\centering\small
\caption{18-month product and business roadmap.}
\begin{tabularx}{\textwidth}{c X l}
\toprule
\textbf{Month} & \textbf{Milestone} & \textbf{Type} \\
\midrule
1--5   & Complete NegoPlay MVP (4-stage pipeline, 4 agents, alignment result)
        & Technical \\
5      & Submit alignment paper to AAMAS / AAAI workshop
        & Academic \\
6      & Open-source profile-discovery code on GitHub
        & Marketing \\
7      & Launch Pro tier on Product Hunt
        & Commercial \\
9      & 20 paying customers; first case study published
        & Commercial \\
12     & HubSpot webhook integration (beta)
        & Technical \\
15     & First enterprise pilot (procurement team)
        & Commercial \\
18     & \$79K ARR; Salesforce connector GA
        & Commercial \\
24     & Series A readiness assessment or academic spin-off decision
        & Strategic \\
\bottomrule
\end{tabularx}
\end{table}

% =====================================================================
% CHAPTER 9 — FINANCIAL PROJECTIONS
% =====================================================================
\chapter{Financial Projections}
\label{ch:finance}

\section{Three-year revenue projection}

\begin{table}[H]
\centering\small
\caption{Three-year revenue, cost, and margin projection.}
\begin{tabularx}{\textwidth}{X r r r}
\toprule
\textbf{Item}                    & \textbf{Year 1} & \textbf{Year 2} & \textbf{Year 3} \\
\midrule
\multicolumn{4}{l}{\textit{Revenue}} \\
Pro seats (avg)                  & 20              & 120             & 600 \\
Pro revenue                      & \$11{,}760      & \$79{,}200      & \$352{,}800 \\
Enterprise / API revenue         & ---             & ---             & \$127{,}200 \\
\textbf{Total ARR}               & \$11{,}760      & \$79{,}200      & \$480{,}000 \\
\midrule
\multicolumn{4}{l}{\textit{Costs}} \\
LLM API -- input tokens  (\$0.075/M) & \$182         & \$1{,}221       & \$7{,}403  \\
LLM API -- output tokens (\$0.300/M) & \$171         & \$1{,}155       & \$6{,}997  \\
LLM API subtotal (Gemini, $\approx 97\,\%$ GM) & \$353         & \$2{,}376       & \$14{,}400 \\
Infrastructure (cloud hosting)   & \$1{,}200       & \$3{,}600       & \$12{,}000 \\
Marketing (LinkedIn, content)    & \$0             & \$24{,}000      & \$60{,}000 \\
\textbf{Total COGS + OpEx}       & \$1{,}553       & \$29{,}976      & \$86{,}400 \\
\midrule
\textbf{Gross Profit}            & \$10{,}207      & \$49{,}224      & \$393{,}600 \\
\textbf{Gross Margin}            & 87\,\%          & 62\,\%          & 82\,\% \\
\bottomrule
\end{tabularx}
\end{table}

\section{Funding requirement}

Phase 1 requires \textbf{zero external capital} (developer is the founder;
LLM costs are under \$15 for the full MVP). Phase 2 marketing costs
(\$24K/year) can be bootstrapped from Phase 1 revenue. A seed round
(\$200K--\$500K) would be sought in Year 2 to fund enterprise sales and
engineering headcount.

% =====================================================================
% CHAPTER 10 — RISKS
% =====================================================================
\chapter{Business Risks}
\label{ch:bizrisks}

\begin{table}[H]
\centering\small
\caption{Commercial risk register.}
\begin{tabularx}{\textwidth}{p{0.5cm} X c c X}
\toprule
\textbf{\#} & \textbf{Risk} & \textbf{Lik.} & \textbf{Imp.} & \textbf{Mitigation} \\
\midrule
BR1 & Research finding not supported ($\rho < 0.5$).
    & Med & High
    & Null result is publishable; pivot to ``bridge coaching AI'' (profile
      analytics without negotiation transfer claim). \\
BR2 & Gong or Salesforce launches a competing feature.
    & Med & High
    & Academic moat (peer-reviewed publication) + niche positioning
      (empirical game-theory basis) differentiates us. \\
BR3 & LLM providers raise prices significantly.
    & Low & Med
    & Provider-neutral \texttt{llm\_client.py} lets us switch to the
      cheapest provider per task within days. \\
BR4 & Data privacy regulations block use of player name data.
    & Low & Med
    & Player data is public (tournament websites); we profile
      \emph{negotiation counterparts}, not users. \\
BR5 & The bridge-to-negotiation transfer claim is rejected by peer reviewers.
    & Med & Med
    & Frame as ``proof of concept in simulated environment'' from the start;
      null result is reported honestly. \\
BR6 & Slow enterprise sales cycle (\textgreater 9 months).
    & High & Med
    & Focus on SME Pro tier first; enterprise is a Year 2+ play. \\
\bottomrule
\end{tabularx}
\end{table}

% =====================================================================
% CHAPTER 11 — TEAM
% =====================================================================
\chapter{Team}
\label{ch:team}

\begin{keybox}[Founder Profile: Anna Ben-Shushan]
\begin{itemize}[leftmargin=*,noitemsep]
    \item \textbf{Current role:} Lecturer at Sami Shamoon College of Engineering;
          PhD candidate at LUT University (Finland).
    \item \textbf{Academic:} M.Sc. Software Engineering; AI Development Expert
          course (current).
    \item \textbf{Industry:} Full-stack developer; AI consultant.
    \item \textbf{PhD topic:} ``From Bridge Tables to Global Markets:
          AI and Optimization Models of Cooperation, Competition, and
          Negotiation on Digital Platforms.''
    \item \textbf{Unique advantage:} Domain expert in both the data source
          (bridge) and the application domain (business negotiation),
          with the academic credibility to publish the alignment finding.
\end{itemize}
\end{keybox}

\section{Hiring plan (if funded)}

\begin{itemize}[leftmargin=*]
    \item \textbf{Month 9:} Part-time front-end developer (student, paid project).
    \item \textbf{Month 15:} Full-time ML engineer (seed-funded).
    \item \textbf{Month 18:} Business development manager for enterprise sales.
\end{itemize}

% =====================================================================
% CHAPTER 12 — CONCLUSION
% =====================================================================
\chapter{Conclusion}
\label{ch:conclusion}

NegoPlay sits at the intersection of three trends: the rapid adoption of
AI in enterprise software, the growing demand for data-driven negotiation
tools, and the academic maturation of multi-agent LLM systems.

Its core differentiator is an \textbf{empirically grounded behavioral
model} -- profiles discovered from 149{,}208 expert decisions rather than
assembled from questionnaires. This gives NegoPlay both a commercial moat
and an academic story that can be defended in peer review.

The financial model is attractive: 97\,\% gross margins on LLM calls,
near-zero CAC in Phase 1, and a clear path to \$480K ARR by Year 3
without external capital for the first 18 months.

\begin{keybox}[The single most important number]
\centering\large
\textbf{Spearman $\rho \geq 0.7$} \\[0.3cm]
\normalsize
If the cross-domain alignment exceeds this threshold, NegoPlay has a
publishable result, a defensible product claim, and a credible reason
for enterprise customers to pay. Everything else follows from this number.
\end{keybox}

\end{document}
